% !TeX program = xelatex
% Overleaf: Menu -> Compiler -> XeLaTeX
% 本文件由 main.tex 精简而来：只保留报告大纲与写作提示，不含现有报告结论或实验数据。
\documentclass[UTF8,11pt,a4paper,fontset=none]{ctexart}

\usepackage{fontspec}
\usepackage{xeCJK}
\IfFontExistsTF{Noto Serif CJK SC}{
  \setCJKmainfont{Noto Serif CJK SC}
  \setCJKsansfont{Noto Sans CJK SC}
  \setCJKmonofont{Noto Sans CJK SC}
}{
  \setCJKmainfont{Songti SC}
  \setCJKsansfont{PingFang SC}
  \setCJKmonofont{PingFang SC}
}

\usepackage[a4paper,margin=2.5cm]{geometry}
\usepackage{amsmath,amssymb,bm}
\usepackage{booktabs,tabularx,array,multirow,longtable}
\usepackage{graphicx,xcolor}
\usepackage{enumitem}
\usepackage{caption,subcaption,float}
\usepackage{microtype}
\usepackage{hyperref}
\usepackage[nameinlink,noabbrev]{cleveref}
\usepackage[numbers,sort&compress]{natbib}

\definecolor{linkblue}{HTML}{175A9C}
\definecolor{promptblue}{HTML}{245A8D}
\hypersetup{
  colorlinks=true,
  linkcolor=linkblue,
  citecolor=linkblue,
  urlcolor=linkblue,
  pdftitle={Bomberman Agent 项目报告},
  pdfauthor={作者姓名}
}

\setlength{\parindent}{2em}
\setlength{\parskip}{0.25em}
\setlist{nosep,leftmargin=2em}
\captionsetup{font=small,labelfont=bf}

% 课程要求：每个章节或小节标题后标明主要作者。
\newcommand{\primaryauthor}[1]{\texorpdfstring{\enspace\normalfont\small[主要作者：#1]}{ [主要作者：#1]}}
\newcommand{\authorsection}[2]{\section{#1\primaryauthor{#2}}}
\newcommand{\authorsubsection}[2]{\subsection{#1\primaryauthor{#2}}}
\newcommand{\authorsubsubsection}[2]{\subsubsection{#1\primaryauthor{#2}}}
\newcommand{\pendingauthor}{待填写}

% 写作提示开关：提交前改为 \draftguidancefalse，所有蓝色提示会消失。
\newif\ifdraftguidance
\draftguidancetrue
\newcommand{\writingnote}[1]{%
  \ifdraftguidance
    \par\noindent\begingroup\color{promptblue}\small
    \fbox{\begin{minipage}{0.94\linewidth}\textbf{写作提示：}#1\end{minipage}}%
    \par\endgroup
  \fi
}
\newcommand{\todoitem}[1]{%
  \ifdraftguidance\textcolor{promptblue}{\textbf{[待补：#1]}}\fi
}

\title{\todoitem{报告标题}}
\author{
  \todoitem{成员姓名 1}\quad
  \todoitem{成员姓名 2}\quad
  \todoitem{成员姓名 3}\\
  \small Machine Learning Essentials, Summer Semester 2026
}
\date{\today}

\begin{document}
\maketitle

\writingnote{提交前确认课程是否接受中文及其字数折算规则。课程要求的总篇幅约为每名成员 4,000 English words，且不应明显超出。报告中不要使用大学 Logo。}

\begin{center}
  \textbf{公开代码仓库：}\todoitem{公开仓库 URL}\\
  \textbf{额外运行依赖：}\todoitem{没有则写“无”；有则列名称与版本}
\end{center}

\begin{abstract}
\writingnote{用一段自洽文字概括：研究问题、比较的模型、统一评估方法、最重要的定量结果、最终模型及选择依据。摘要最后写清结论的适用边界；不要引用文献，也不要加入正文未报告的结果。}
\todoitem{摘要正文}
\end{abstract}

\noindent\textbf{关键词：}\todoitem{3--5 个关键词}

\tableofcontents
\clearpage

\writingnote{全文必须覆盖团队开发的所有模型，并说明最终模型的选择依据。建议把最多篇幅留给“实验与结果”；扩写应来自实验假设、协议、结果、失败解释和决策依据，而不是重复算法定义。}

\authorsection{引言}{\pendingauthor}
\writingnote{简要说明问题、研究价值和主要挑战；提出全文研究问题并概括贡献。避免在此堆叠所有实验数字。}

\authorsubsection{问题与研究动机}{\pendingauthor}
\writingnote{说明 Bomberman 中导航、稀疏奖励、延迟爆炸、多智能体竞争与 CPU 推理时限为何构成机器学习问题。}
\todoitem{正文}

\authorsubsection{研究问题与成功标准}{\pendingauthor}
\writingnote{列出 2--4 个可由实验回答的研究问题。为每个问题定义可观测指标和判断标准，例如得分、金币、击杀、生存率、全收集率、WAIT/循环和完整 act 延迟。}
\todoitem{正文}

\authorsubsection{贡献与全文结构}{\pendingauthor}
\writingnote{只概括有证据支持的贡献，并用一小段介绍后续章节。不要把计划或未验证功能写成贡献。}
\todoitem{正文}

\authorsection{背景}{\pendingauthor}
\writingnote{比引言更具体地描述任务，并概述团队考虑过的强化学习方法。引用原始或权威文献，避免写成通用教材摘录。}

\authorsubsection{Bomberman 任务与环境}{\pendingauthor}
\writingnote{交代状态、六个动作、奖励目标、炸弹时序、障碍、回合长度和正式 CPU/内存/决策时间约束。所有数值应回查课程材料与源码。}
\todoitem{正文}

\authorsubsection{Task 1--4 课程任务}{\pendingauthor}
\writingnote{解释从金币导航、炸箱逃生、弱对手到强对手的难度递进，以及每阶段测试的能力。区分课程建议和团队自定门槛。}
\todoitem{正文}

\authorsubsection{强化学习形式化与候选方法}{\pendingauthor}
\writingnote{定义状态、动作、奖励、策略和价值函数；概述实际考虑的 Q-learning、DQN、Double DQN、n-step、PER、蒸馏或其他方法。只写与项目决策相关的理论。}
\todoitem{正文与必要公式}

\authorsubsection{奖励塑形、课程学习与安全约束}{\pendingauthor}
\writingnote{说明这些概念为什么适合本任务、可能引入什么偏差，以及之后将如何用消融或冻结评估验证。}
\todoitem{正文}

\authorsection{项目规划}{\pendingauthor}
\writingnote{说明真实的团队组织、时间规划、子任务分配和协作方式。不能把项目描述成成员各自独立完成一个模型；应展示共同设计、评审、工具或实验决策。}

\authorsubsection{团队组织与共同目标}{\pendingauthor}
\writingnote{说明成员角色、共同会议/评审方式、共享代码和实验协议，以及关键决定如何形成。}
\todoitem{正文}

\authorsubsection{分工与共享成果}{\pendingauthor}
\writingnote{用表格列责任人、工作包、协作者、输出和证据路径。明确哪些特征、奖励、安全层或评估工具被多条模型路线共同使用。}
\todoitem{正文或分工表}

\authorsubsection{里程碑、晋级与停止规则}{\pendingauthor}
\writingnote{给出真实时间线和 Task 1--4 的进入/退出条件。停止某条路线时写清门槛、证据、资源权衡与信息增益，不要只说“效果不好”。}
\todoitem{正文或时间线}

\authorsubsection{计算资源与证据管理}{\pendingauthor}
\writingnote{记录训练硬件、预算、随机种子、checkpoint、run ID、逐局输出、证据状态和复现路径。使用深度学习时必须说明训练硬件资源。}
\todoitem{正文}

\authorsection{方法}{\pendingauthor}
\writingnote{解释并论证关键设计选择、待比较变体、性能指标和系统测试方法，并引用相关文献。方法章写“怎么做、为什么”，结果章写“发生了什么”。}

\authorsubsection{状态表示与特征}{\pendingauthor}
\writingnote{逐模型说明输入形状、归一化、历史信息、空间通道或离散编码。解释每组特征对应的失败假设，避免把版本号直接当成性能顺序。}
\todoitem{正文与特征表}

\authorsubsection{模型与学习目标}{\pendingauthor}
\writingnote{团队必须开发并比较至少两个不同模型。对每个模型说明结构、输出、损失/更新公式、目标网络或估计器、关键超参数、相对基线的变化和预期风险。}
\todoitem{正文、公式与架构图}

\authorsubsection{奖励设计}{\pendingauthor}
\writingnote{区分正式游戏得分和训练辅助奖励；给出数值、触发条件、信用归因及防止刷分/循环的设计。多因素版本不能被描述成单因素因果结论。}
\todoitem{正文与奖励表}

\authorsubsection{动作合法性与安全机制}{\pendingauthor}
\writingnote{区分物理合法、危险预测、动作过滤和学习器排序。说明 fallback、规划时域、计算代价与已知反例；若使用规则层，论证它没有直接替代模型决定“最佳动作”。}
\todoitem{正文}

\authorsubsection{探索、迁移与推理流程}{\pendingauthor}
\writingnote{描述探索策略、replay、target 更新、checkpoint resume、policy-only transfer 和推理阶段。给出一张从 game state 到动作输出的流程图会更清楚。}
\todoitem{正文}

\authorsubsection{评估指标与统计协议}{\pendingauthor}
\writingnote{预先定义不同 Task 的主/次指标、基线、对手、场景、局数、环境种子、独立训练种子、冻结条件和统计方法。说明何时采用配对比较、置信区间或 bootstrap。}
\todoitem{正文与协议表}

\authorsection{训练}{\pendingauthor}
\writingnote{描述训练过程和用于加速或改善训练的技巧。按阶段说明初始化、训练预算、checkpoint 选择、迁移方式、旧任务保持和失败纠正。}

\authorsubsection{分阶段训练课程}{\pendingauthor}
\writingnote{逐 Task 说明场景、对手、允许动作、训练目标、晋级门槛和实际到达阶段。未运行的阶段明确标为未运行，不能填成零。}
\todoitem{正文或课程表}

\authorsubsection{超参数与训练配置}{\pendingauthor}
\writingnote{提供足以复现的参数表：优化器、学习率、折扣率、batch、replay、探索计划、目标同步、训练局数/步数、种子与设备。注明模型间哪些参数不同。}
\todoitem{配置表}

\authorsubsection{Checkpoint、Resume 与 Transfer}{\pendingauthor}
\writingnote{解释 checkpoint 保存内容、严格恢复与仅策略迁移的区别，以及特征/奖励/安全合同变化时如何避免误加载旧状态。}
\todoitem{正文}

\authorsubsection{训练效率与加速技巧}{\pendingauthor}
\writingnote{报告无 GUI、并行训练、GPU、缓存或数据复用等做法及其实际收益。最终提交 Agent 不得使用 multiprocessing；训练可使用。}
\todoitem{正文与测量结果}

\authorsubsection{训练失败与纠正}{\pendingauthor}
\writingnote{保留未收敛、遗忘、循环、自杀、指标缺陷或实现错误。说明诊断证据、修改内容，以及修改是否改变策略本身或只修复测量。}
\todoitem{正文}

\authorsection{实验与结果}{\pendingauthor}
\label{sec:results}
\writingnote{本章是报告重点。每组实验按“问题/假设—协议—结果—解释—决策”书写。用冻结策略和正式游戏指标，不用训练奖励代替评估。必须报告没有改善或失败的结果。}

\authorsubsection{总体实验逻辑与比较边界}{\pendingauthor}
\writingnote{说明哪些实验可直接比较，哪些因场景、对手、seed、checkpoint 或统计单位不同只能诊断。训练 seed 和同一模型的环境重复不能混为独立训练。}
\todoitem{正文}

\authorsubsection{基线与 Task 1：导航能力}{\pendingauthor}
\writingnote{比较课程基线和各学习模型。报告金币、全收集率、步数、非法动作、WAIT/循环及不确定性；说明晋级决定。}
\todoitem{协议、表/图、结果解释与决定}

\authorsubsection{Task 2：炸箱、逃生与目标重选}{\pendingauthor}
\writingnote{报告金币、炸箱、炸弹存活、自杀、全收集、停滞和旧任务保持。把安全改善与长期规划改善分开讨论。}
\todoitem{协议、表/图、结果解释与决定}

\authorsubsection{Task 3：弱对手下的竞争}{\pendingauthor}
\writingnote{报告得分、金币、击杀、名次、生存及旧任务回归。若比较父子 checkpoint，优先在相同世界/seed 上做配对分析。}
\todoitem{协议、表/图、结果解释与决定}

\authorsubsection{Task 4：强对手与最终候选}{\pendingauthor}
\writingnote{区分开发筛选、独立确认、历史 benchmark 和专项测试。报告策略稳定性、失败模式和计算成本；不要把不同协议合并成统一排行榜。}
\todoitem{协议、表/图、结果解释与决定}

\authorsubsection{消融、失败结果与误差分析}{\pendingauthor}
\writingnote{至少覆盖关键特征、奖励、安全机制或训练技巧的对照。展示代表性失败轨迹，并判断瓶颈来自表示、奖励、学习、动作过滤还是评估。}
\todoitem{消融表、失败案例与解释}

\authorsubsection{最终跨模型比较与选择}{\pendingauthor}
\writingnote{在完全一致的场景、对手、seed、局数和冻结条件下比较最终候选与课程基线。说明主指标、置信区间、旧任务保持、CPU 完整 act 延迟、内存、超时和打包等价性。}
\todoitem{最终比较表与选择依据}

\authorsection{结论}{\pendingauthor}
\writingnote{总结发现并直接回答研究问题；解释最终选择的证据与适用边界。不要在结论引入新实验。}

\authorsubsection{主要发现与研究问题回答}{\pendingauthor}
\writingnote{逐条对应引言中的研究问题，给出最关键的定量证据和最终模型选择。}
\todoitem{正文}

\authorsubsection{局限}{\pendingauthor}
\writingnote{讨论样本量、训练 seed、协议缺口、计算资源、未覆盖对手、测量范围、语言/篇幅确认和无法证明的安全性质。}
\todoitem{正文}

\authorsubsection{未来工作}{\pendingauthor}
\writingnote{说明如果有更多时间，最值得执行的 Agent 改进、实验和游戏设置变化；把建议与已观察失败对应起来。}
\todoitem{正文}

\bibliographystyle{plainnat}
\bibliography{references}

\appendix
\authorsection{模型清单与证据索引}{\pendingauthor}
\writingnote{列出所有开发模型、配置、checkpoint、哈希、训练终点、证据状态和路径。附录不能替代正文中的核心方法与结果。}
\todoitem{清单表}

\authorsection{补充实验与复现信息}{\pendingauthor}
\writingnote{放置正文容纳不下但复现所需的命令、完整超参数、额外结果、数据契约和打包检查。注意最终公开仓库不能包含报告本身。}
\todoitem{补充内容}

\end{document}
